%!TEX root = ../main.tex
% =============================================================
%  CAPÍTULO 1 — INTRODUCCIÓN            (Entrega 1 · 5 % · 9/6)
% =============================================================

\chapter{Introducción}\label{cap:introduccion}

\begin{guia}[Guía de la cátedra: qué se evalúa en este capítulo]
\begin{enumerate}
    \item \textbf{Descripción del problema}: concreto, anclado en evidencia
          (datos duros con fuente, casos reales, competidores nombrados).
    \item \textbf{Objetivo general}: uno solo, medible, en un párrafo.
          Los \textbf{objetivos específicos} se desprenden de él y
          \emph{no} son los pasos del plan de trabajo.
    \item \textbf{Hipótesis}: explícita, formulada como pregunta o
          afirmación de investigación que pueda responderse en el
          capítulo~\ref{cap:conclusiones}.
    \item \textbf{Relevancia}: por qué importa, a quién beneficia, qué les
          falta a las soluciones existentes.
    \item \textbf{Alcance y recursos}: qué entra y qué no; stack y fuentes
          de datos previstos. El ``alcance que nunca cierra'' es el riesgo
          número uno de la materia.
\end{enumerate}
Extensión orientativa: 4 a 6 páginas. Redacción impersonal, sin ``yo
creo'', oraciones de 15 a 22 palabras, sin palabras vacías
(``claramente'', ``obviamente''). Toda cifra con su fuente citada, por
ejemplo \citep{breiman2001}.
\end{guia}

La organización de un evento social o empresarial requiere coordinar una
gran cantidad de servicios, como catering, fotografía, decoración,
ambientación y entretenimiento, entre otros. Generalmente estos servicios
se contratan por separado, por lo que la persona que organiza el evento
debe utilizar distintas plataformas, redes sociales y medios de
comunicación para buscar y contratar cada uno de ellos. En muchos casos
también es necesario contratar a un organizador o coordinador que actúe
como intermediario entre todos los proveedores y ayude en la toma de
decisiones relacionadas con la planificación del evento.

Si bien hoy existen plataformas para buscar proveedores y otras para
administrar reservas de salones, la mayoría de estas soluciones funcionan
por separado y solo resuelven una parte del problema. Ninguna permite
gestionar todo el evento en un mismo lugar ni brinda herramientas que
ayuden a anticipar situaciones, evaluar distintas alternativas y facilitar
la toma de decisiones. Ante esta situación, surge la necesidad de
desarrollar una plataforma que centralice la contratación de servicios, la
gestión de salones y proveedores y la organización integral del evento
dentro de un único entorno digital.

\section{Pregunta de investigación}

\completar{¿En qué medida la integración, dentro de una misma plataforma,
de un marketplace de servicios, un sistema de recomendación basado en
\gls{ia} y un Gemelo Digital del evento mejora la eficiencia del proceso
de organización de eventos y la calidad de las decisiones tomadas por
clientes, salones y proveedores, en comparación con el uso de herramientas
fragmentadas e independientes como el que se releva actualmente en la
industria?}

\section{Objetivo general}

\begin{ejemplo}[Ejemplo: forma de un objetivo general]
Desarrollar un sistema de \completar{qué} para \completar{quién} en
\completar{dónde}, que permita \completar{resultado medible} mediante
\completar{enfoque o tecnología central}.
\end{ejemplo}

Determinar en qué medida la integración de un marketplace de servicios, un
sistema de recomendación basado en \gls{ia} y un Gemelo Digital del
evento, dentro de una misma plataforma, mejora la eficiencia del proceso
de organización de eventos y la calidad de las decisiones tomadas por los
distintos participantes de la industria (clientes, salones y
proveedores), frente al uso de herramientas fragmentadas.

\subsection{Propósito del proyecto (producto)}

Para responder la pregunta de investigación, este proyecto desarrolla una
plataforma inteligente que centraliza la contratación de servicios para
eventos, la gestión comercial y operativa de salones y proveedores, y la
planificación del evento mediante recomendación y simulación con Gemelos
Digitales.

El objetivo de investigación y el producto son cosas distintas: el
objetivo es lo que la tesis busca averiguar, mientras que la plataforma
es la herramienta que se construye para poder medirlo.

\section{Objetivos específicos}

\begin{itemize}
    \item Diseñar una plataforma que permita a los usuarios buscar,
          comparar y contratar los diferentes servicios necesarios para la
          organización de un evento desde un mismo lugar.
    \item Desarrollar herramientas de gestión para que los salones y
          proveedores puedan administrar sus clientes, reservas,
          presupuestos, contratos y pagos de manera más organizada.
    \item Implementar funcionalidades que permitan gestionar la
          disponibilidad de los servicios y recursos involucrados en cada
          evento.
    \item Incorporar herramientas inteligentes que ayuden a recomendar
          servicios y alternativas de organización de acuerdo con las
          características y necesidades de cada usuario.
    \item Desarrollar un Gemelo Digital que permita representar el evento
          de forma virtual y analizar distintos escenarios antes de su
          realización.
    \item Generar información e indicadores que faciliten la toma de
          decisiones tanto para los clientes como para los salones y
          proveedores involucrados en la organización del evento.
\end{itemize}

\section{Barreras}

La principal dificultad para desarrollar este proyecto es poder reunir en
una misma plataforma todas las necesidades de los distintos participantes
de la industria de eventos, como los clientes, los salones y los
proveedores, ya que cada uno maneja información y procesos diferentes.
Además, otro desafío es desarrollar herramientas inteligentes que sean
capaces de recomendar servicios y simular distintos escenarios del evento,
para lo cual es necesario contar con información que permita representar
de la mejor manera posible cómo se comporta la organización de un evento.
Por último, también representa un desafío coordinar la disponibilidad de
múltiples servicios y proveedores al mismo tiempo, y desarrollar una
plataforma que pueda manejar una gran cantidad de usuarios e información
sin perder rendimiento.

\section{Relevancia}

La relevancia de este proyecto se encuentra en la posibilidad de reunir en
un solo lugar todo el proceso de organización de un evento, facilitando la
coordinación entre los distintos participantes y simplificando la toma de
decisiones.

La plataforma busca beneficiar tanto a los clientes, que pueden organizar
y contratar todos los servicios desde un mismo lugar, como a los salones y
proveedores, que disponen de herramientas para gestionar de forma más
ordenada sus eventos, clientes y contrataciones.

Además, la incorporación de herramientas de simulación y recomendación le
agrega un componente innovador al proyecto, ya que la plataforma no solo
permite contratar servicios, sino también ayudar a planificar el evento y
evaluar diferentes alternativas antes de tomar una decisión.

\section{Hipótesis}

Se considera que reunir en una única plataforma la contratación de
servicios, la gestión de salones y proveedores, y herramientas
inteligentes de recomendación y simulación permite hacer más simple y
eficiente la organización de eventos, mejorando la coordinación entre los
distintos participantes y facilitando la toma de decisiones durante todo
el proceso de planificación. Asimismo, la incorporación de un Gemelo
Digital del evento permite representar de manera virtual diferentes
escenarios antes de su realización, anticipando posibles situaciones y
evaluando el impacto de las decisiones tomadas sobre la organización
integral del evento.

\section{Enfoque}

El proyecto se desarrolla en tres etapas incrementales: primero el
relevamiento de los procesos actuales y la construcción de la
arquitectura, el marketplace y las herramientas de administración; luego
la incorporación de la recomendación de servicios y la simulación de
escenarios del evento; y finalmente las pruebas y la evaluación del
funcionamiento y el aporte de la plataforma.

\section{Recursos preliminares}

Para la realización del proyecto se cuenta con la participación del
alumno como responsable del análisis, diseño, desarrollo y validación de
la plataforma. También se prevé la colaboración de salones de eventos,
proveedores y organizadores, quienes aportan información sobre sus
procesos de trabajo y permiten evaluar el funcionamiento de la solución
propuesta.

En cuanto a los recursos tecnológicos, se utilizan herramientas de
desarrollo web, una base de datos y tecnologías que permiten implementar
las funcionalidades del marketplace, la gestión de salones y proveedores y
las herramientas inteligentes de recomendación y simulación. Además, se
utilizan herramientas de generación de reportes e indicadores que
permiten visualizar información relevante y analizar diferentes
escenarios de los eventos mediante Gemelos Digitales.
